\documentclass[11pt]{article}

\usepackage{epsfig}
\usepackage{amsfonts}
\usepackage{amssymb}
\usepackage{amstext}
\usepackage{amsmath}
\usepackage{xspace}
\usepackage{theorem}
\usepackage{hyperref}
\usepackage{fullpage}
%\usepackage[]{algorithm2e}
\usepackage{xfrac}
\usepackage{mathtools}
\usepackage{algorithm}
\usepackage{algpseudocode}

\usepackage{listings}

\DeclarePairedDelimiter\ceil{\lceil}{\rceil}
\DeclarePairedDelimiter\floor{\lfloor}{\rfloor}

\usepackage{enumitem}                     

\usepackage{tikz}
\usepackage{tikz-qtree}
\usetikzlibrary{shapes}

\tikzstyle{code} = [black!90, draw=black!30, fill=black!5, very thick,
    rectangle, dashed, inner xsep=10pt, inner ysep=7pt]

\newenvironment{codebox}{
    \hspace{.05\textwidth}
        \begin{tikzpicture}
            \node[code] \bgroup
                \begin{minipage}{.65\textwidth}
                    \begin{alltt}}
                    {\end{alltt}
                \end{minipage}
            \egroup;
        \end{tikzpicture}
}

\newif\ifrubric
\rubrictrue
\rubricfalse

% This is the stuff for normal spacing
%\makeatletter
% \setlength{\textwidth}{6.5in}
% \setlength{\oddsidemargin}{0in}
% \setlength{\evensidemargin}{0in}
% \setlength{\topmargin}{0.25in}
% \setlength{\textheight}{8.25in}
% \setlength{\headheight}{0pt}
% \setlength{\headsep}{0pt}
% \setlength{\marginparwidth}{59pt}
%
% \setlength{\parindent}{0pt}
% \setlength{\parskip}{5pt plus 1pt}
% \setlength{\theorempreskipamount}{5pt plus 1pt}
% \setlength{\theorempostskipamount}{0pt}
% \setlength{\abovedisplayskip}{8pt plus 3pt minus 6pt}
 
 
 \usepackage{titlesec}

\titleformat*{\section}{\bfseries}
\titleformat*{\subsection}{\bfseries}
\titleformat*{\subsubsection}{\bfseries}
\titleformat*{\paragraph}{\bfseries}
\titleformat*{\subparagraph}{\bfseries}

% \renewcommand{\section}{\@startsection{section}{1}{0mm}%
%                                   {2ex plus -1ex minus -.2ex}%
%                                   {1.3ex plus .2ex}%
%                                   {\normalfont\Large\bfseries}}%
% \renewcommand{\subsection}{\@startsection{subsection}{2}{0mm}%
%                                     {1ex plus -1ex minus -.2ex}%
%                                     {1ex plus .2ex}%
%                                     {\normalfont\large\bfseries}}%
% \renewcommand{\subsubsection}{\@startsection{subsubsection}{3}{0mm}%
%                                     {1ex plus -1ex minus -.2ex}%
%                                     {1ex plus .2ex}%
%                                     {\normalfont\normalsize\bfseries}}
% \renewcommand\paragraph{\@startsection{paragraph}{4}{0mm}%
%                                    {1ex \@plus1ex \@minus.2ex}%
%                                    {-1em}%
%                                    {\normalfont\normalsize\bfseries}}
% \renewcommand\subparagraph{\@startsection{subparagraph}{5}{\parindent}%
%                                       {2.0ex \@plus1ex \@minus .2ex}%
%                                       {-1em}%
%                                      {\normalfont\normalsize\bfseries}}
%\makeatother

\newenvironment{proof}{{\bf Proof:  }}{\hfill\rule{2mm}{2mm}}
\newenvironment{proofof}[1]{{\bf Proof of #1:  }}{\hfill\rule{2mm}{2mm}}
\newenvironment{proofofnobox}[1]{{\bf#1:  }}{}
\newenvironment{example}{{\bf Example:  }}{\hfill\rule{2mm}{2mm}}
%\renewcommand{\thesection}{\lecnum.\arabic{section}}

%\renewcommand{\theequation}{\thesection.\arabic{equation}}
%\renewcommand{\thefigure}{\thesection.\arabic{figure}}

%\renewcommand{\theequation}{\lecnum.\arabic{equation}}
%\renewcommand{\thefigure}{\lecnum.\arabic{figure}}

%\newcounter{LecNum}
%\setcounter{LecNum}{1}

%\newtheorem{fact}{Fact}[LecNum]
\newtheorem{fact}{Fact}
\newtheorem{lemma}[fact]{Lemma}
\newtheorem{theorem}[fact]{Theorem}
\newtheorem{definition}[fact]{Definition}
\newtheorem{corollary}[fact]{Corollary}
\newtheorem{proposition}[fact]{Proposition}
\newtheorem{claim}[fact]{Claim}
\newtheorem{exercise}[fact]{Exercise}

% math notation
\newcommand{\R}{\ensuremath{\mathbb R}}
\newcommand{\Z}{\ensuremath{\mathbb Z}}
\newcommand{\N}{\ensuremath{\mathbb N}}
\newcommand{\F}{\ensuremath{\mathcal F}}
\newcommand{\SymGrp}{\ensuremath{\mathfrak S}}

\newcommand{\size}[1]{\ensuremath{\left|#1\right|}}
%\newcommand{\ceil}[1]{\ensuremath{\left\lceil#1\right\rceil}}
%\newcommand{\floor}[1]{\ensuremath{\left\lfloor#1\right\rfloor}}
\newcommand{\poly}{\operatorname{poly}}
\newcommand{\polylog}{\operatorname{polylog}}

% anupam's abbreviations
\newcommand{\e}{\epsilon}
\newcommand{\half}{\ensuremath{\frac{1}{2}}}
\newcommand{\junk}[1]{}
\newcommand{\sse}{\subseteq}
\newcommand{\union}{\cup}
\newcommand{\meet}{\wedge}

\newcommand{\prob}[1]{\ensuremath{\text{{\bf Pr}$\left[#1\right]$}}}
\newcommand{\expct}[1]{\ensuremath{\text{{\bf E}$\left[#1\right]$}}}
\newcommand{\Event}{{\mathcal E}}
\newcommand{\E}{\ensuremath{\text{E}}}

\newcommand{\mnote}[1]{\normalmarginpar \marginpar{\tiny #1}}

\setenumerate[0]{label=(\alph*)}


%%%%%%%%%%%%%%%%%%%%%%%%%%%%%%%%%%%%%%%%%%%%%%%%%%%%%%%%%%%%%%%%%%%%%%%%%%%
% Document begins here %%%%%%%%%%%%%%%%%%%%%%%%%%%%%%%%%%%%%%%%%%%%%%%%%%%%
%%%%%%%%%%%%%%%%%%%%%%%%%%%%%%%%%%%%%%%%%%%%%%%%%%%%%%%%%%%%%%%%%%%%%%%%%%%



\begin{document}

\noindent {\large {\bf 601.433/633 Introduction to Algorithms} \hfill {{\bf Fall 2026}}}\\
{{\bf Homework \#3}} \hfill {{\bf Due:} October 1, 2026, 11:59pm} \\
\rule[0.1in]{\textwidth}{0.4pt}

Remember: you may work in groups, but must write up your solution entirely on your own.  Solutions must by typeset (\LaTeX\ preferred but not required). 
Collaboration is limited to discussing the problems -- you may not look at, compare, reuse, etc.~any text from anyone else in the class. 
You may use the internet to look up formulas, definitions, etc., but may not simply look up the answers online.  

Please include proofs with all of your answers, unless stated otherwise.

\noindent \rule[0.1in]{\textwidth}{0.4pt}

\section{Amortized analysis of 2-3-4 trees (33 points)}

Recall that in a 2-3-4 tree, whenever we insert a new key we immediately split (on our way down the tree) any node we see that is full (has 3 keys in it).  

\begin{enumerate}
\item (11 points) Show that in the worst case, the number of splits that we do in a single operation can be $\Omega(\log n)$. 
In other words, show that there is a series of $n$ inserts so that the next insert causes $\Omega(\log n)$ splits. 
You can be informal here --- just explain at a high level what such a sequence would look like and why it would result in $\Omega(\log n)$ splits.

The key observation in 2-3-4 trees is that every leaf must have the same depth. 
Our worst case scenario is that every node along the insertion path is full (a 4-node), which means that on the way down, every node must be split. 
Therefore, $1$ split occurs for every level of the tree. To find the lower bound for the number of splits, we must find the depth of the tree. 
Because each level increases the number of nodes by a constant level (between 2-4, since the number of children of any internal node is between 2 and 4), 
we can represent the number of levels in the tree as $\Theta(\log n)$ (as there are currently $n$ nodes in the tree). 
To insert a new element, we perform $1$ split at each of the $\Theta(\log n)$ levels, yielding $\Omega(\log n)$ splits.

\end{enumerate}

Let's try to get around this worst-case bound by using amortized analysis.  We will try to prove that the amortized number of splits is only $O(1)$. 
Note that this does not mean that the amortized running time of an insert is $O(1)$ (since this is not true); 
it just means that the amortized number of splits is $O(1)$.  So think of ``cost'' not as ``time'', but as ``number of splits''. 
Since we didn't talk about them in class, feel free to assume there are no delete operations.  

\begin{enumerate}[resume]
\item (11 points) Since we only have to split a node when it's full, a natural approach is to have a bank on every node which works as follows. 
When a node is first created its bank is initialized to $0$.  When a node becomes full, we add $1$ to its bank. 
Then when we split a node we use the $1$ in its bank to pay for the split.  In other words, a node's bank is $1$ if the node is full and $0$ otherwise.

This argument unfortunately does not work.  Explain why this banking scheme \emph{does not} imply that the amortized number of splits is $O(1)$. 
Your explanation should be convincing, but does not need to be a fully formal proof

This approach is trying to prove that over $m$ insertions, the number of splits is bounded by $O(m)$. 
We must think of the case above, where the insertion triggers $\Theta(\log n)$ splits. 
Using this approach, this means that each of those $\Theta(\log n)$ nodes has $1$ token stored, which are then used to pay for the split. 
Thus, the banking scheme does accurately show that each split has the tokens to pay for it.

However, the approach does not establish that the amount added to the bank over the sequence is sufficient to bound the total number of split. 
Simply showing that each split has the tokens to pay for it does not establish that the total number of splits is bounded by $O(m)$, 
as the tokens used to pay for the split may have been accumulated over many previous insertions. 
This bank approach does show that each individual $m$ split can be paid for, but it does not show that the number of splits is bounded by $O(m)$
Therefore, although the bank can pay for each individual split, this does not prove a constant amortized number of splits per insertion.

\item (11 points) Show how to modify this argument to work: use a bank at each node to prove that the amortized number of splits is $O(1)$. 
Hint: the previous part shows that this bank cannot just equal $1$ if the node is full and $0$ otherwise.

I will use the approach that each time an element is added to one of the nodes in the tree $1$ token is added to that node's bank. 
Thus, since a node is full when it has $3$ elements, it will have accumulated $3$ tokens. 
We can establish that over the course of $m$ insertions, $m$ tokens have been added to the bank. 

Now, we can move on to proving that the total number of splits is bounded by $O(m)$. 
On the splitting of a node, $1$ token is spent, leaving $2$ tokens. 
Since a split occurs when a node has $3$ elements, the resulting node has $2$ children. 
Thus, each of those $2$ resulting nodes of $1$ element each can get $1$ token. 

Therefore, we know that every split that occurs can be paid for by a token that has been previously deposited. 
Since each of the $m$ insertions deposits $1$ token, at most $m$ tokens are deposited total. 
Therefore, we can prove that there are at most $m$ splits over $m$ insertions, which creates the upper bound $O(m)$ for the number of splits. 
Since we have shown $O(m)$ splits over $m$ insertions, we can therefore conclude a constant amortized number of splits per insertion.

TODO

\end{enumerate}

\section{Move to Front (34 points)}
Amortized analysis can be used not just to give absolute bounds on running times, but also bounds \emph{compared to other algorithms}. 
In this problem we'll see an example of this.

Suppose we have $n$ items $x_1, x_2, \dots, x_n$ that we wish to store in a linked list.  The cost of a \emph{lookup} operation is the position in the list, 
i.e.~if we lookup some element $x_i$ we pay $1$ if it is the first element, $2$ if it is the second element, $3$ if it is the third element, etc. 
This models the time it takes to scan through the list to find the element.  

For example, suppose there are four items and we lookup $x_1$ twice, $x_2$ three times, $x_3$ once, and $x_4$ four times. 
Then if we store them in the list $(x_1, x_2, x_3, x_4)$ the total cost is $1(2) + 2(3) + 3(1) + 4(4) = 27$. 
On the other hand, if we stored them in the list $(x_4, x_2, x_1, x_3)$ the cost would be $1(4) + 2(3) + 3(2) + 4(1) = 20$. 
It is easy to see that if we knew the number of times each element was looked up, the optimal list is simply sorting by the number of lookups, 
i.e.~the first element is the one looked up most often, then the element looked up next most often, etc.  

But what if we do not know how many times each element will be looked up?  It turns out that a good strategy is \emph{Move-to-Front} (MTF): 
when we lookup an item, we also move it to the head of the list.  Let's say that moving an element is free 
(it is easy enough to implement with a constant number of pointer switches, in any case). 
So if, for example, we start with the list $(x_1, x_2, \dots, x_n)$ and then lookup $x_4$, we pay a cost of $4$ and 
afterwards the list will be $(x_4, x_1, x_2, x_3, x_5, x_6, \dots, x_n)$.  Then if we lookup $x_4$ again we only have to pay a cost of $1$.  

\begin{enumerate}

\item (17 points) Consider a sequence of $m$ operations (think of $m$ as being much larger than $n$). 
Let $C_{init}$ be the total cost of these operations if we used the original list $(x_1, x_2, \dots, x_n)$ the entire time (i.e., we do not use MTF). 
Let $C_{MTF}$ be the total cost if we use MTF on the same operations starting from the original list.  Prove that $C_{MTF} \leq 2 C_{init}$.  

Hint: Think of each item $x_i$ as having a ``piggy bank'' in which the number of tokens is the number of elements before it in the current list 
that come after it in the initial list, or a potential function equal to the number of reversals.  

Let us assume that for each element $x_i$, where $i$ is the index of the element in the list. 
This means that the cost $C_i$ of accessing that element is $i$. 
We will consider a sequence of $m$ lookup operations, and consider the cost $C$ of this series of operations. 

Let us consider the cost of looking up some element $x_i$. 
In the initial lookup implementation, the cost would be: $C_i = i$

However, let us consider this cost of access for element $x_i$ with the MTF lookup operation. 
In this, we must consider both the actual access cost and the tokens accumulated from previous reversals. 
For any the access of element $x_i$, we can consider $1$ token stored for each element that is currently before $x_i$ which was initially after $x_i$, 
since the process of moving $x_i$ to the front reverses back what was already reversed. 

Similar to the naive implementation, the initial lookup cost is: $C_i = i$. 
Since this applies to any element, we can generalize this to any lookup with the initial algorithm: $C_{init} = i$.

After a series of some operations, $x_i$ will have changed location. We set this location to $j$. 
We can divide the set of elements before $x_i$ after some operations into two categories: 
elements that were before $x_i$ in the initial list, which we term $B$ and elements that were originally after $x_i$, which we term $D$. 
Therefore, we can represent the current position $j$ as: $j = |B| + |D| + 1$

We can represent the new cost of access of $x_i$ as: $C_i = |B| + |D| + 1$. 
We know that every element $D$ which is now before $x_i$ started out after $x_i$, so there are $|D|$ tokens in the $x_i$'s piggy bank. 
We can spend these tokens to cover $|D|$ units of the actual access cost, leaving $|B| + 1$ units to be paid. 

As we move $x_i$ to the front, $|B|$ tokens enter $x_i$'s piggy bank, since they represent elements that were before and are now after $x_i$. 
However, we must still cover the cost $|B| + 1$. Because $x_i$ was initially at position $i$, we know that $|B| \leq (i-1)$. 
Therefore, the maximum actual cost of any access of $x_i$ that is left to be paid after paying the accumulated $|D|$ tokens is: $|B| + 1 \leq i$.

In the process of accessing $x_i$, we must create the $|B|$ tokens, while also actually accessing the element. 
Thus, the operation must access the element, spend the $|D|$ tokens, and accumulate the $|B|$ tokens. 
We can express the accounting cost $A_i$ for the number of tokens as: 

$A_i =  C_i - \text{(tokens spent)} + \text{(tokens created)}$  

$A_i = |B| + |D| + 1 - |D| + |B|$. 


This simplifies to: $A_i = 2|B| + 1$. 

Using the bound we found for $|B|$, we get: 
$A_i \leq 2(i - 1) + 1 = 2i - 1 < 2i$

Therefore, the accounting cost of this operation is at most $2i$. 
The sum of the accounting costs is an upper bound for actual cost of the sequence of MTF operations. 
This is because each time a token is spent, it was deposited earlier, and each time a token is deposited, it contributes to the accounting cost. 
Therefore, this $2i$ upper bound represents the tokens that were deposited but never needed. Because there are initially zero tokens 
and every token spent was deposited, the sum of the accounting cost is at least the total MTF cost. 
From here, we sum the accounting cost bounds over the sequence of $m$ operations.

$C_{MTF} \leq (2i_1 + 2i_2 + 2i_3 + ... + 2i_m)$

Now summing the initial algorithm, since the list nevver changes:

$C_{init} = i_1 + i_2 + i_3 + ... + i_m$

Putting this together:

$C_{MTF} \leq (2i_1 + 2i_2 + 2i_3 + ... + 2i_m) = 2 \cdot C_{init}$

Therefore:

$C_{MTF} \leq 2 \cdot C_{init}$

\item (17 points) Now let $C_{static}$ be the smallest possible cost among all fixed lists.  In other words, given the same sequence of lookups, 
each fixed list (without MTF) has some cost, and let $C_{static}$ be the minimum of those costs. 
Note that this will correspond to the list that is sorted by the number of accesses, like in the first example.  Prove that $C_{MTF} \leq 2C_{static} + n^2$.

Hint: Think about your bank or potential function argument from the previous part. 
What if you redefine it to be relative to the optimal list, rather than the initial list?  Then what is the initial bank/potential value? 
How does this affect the amortized analysis?

We define $L$ to be the list where the cost of the $m$ operations is $C_{static}$.

In the last problem, we compared the MTF list to the original list. In this list, we define it relative to the optimal list,  organized by number of accesses. 
I define the cost of each access to element $x_i$ in the MTF algorithm similar to in the previous part; 
however, instead of considering sets $F$ and $G$ to contain the elements originally before and after $x_i$, 
we define $F$ to be the elements that exist in $L$ before $x_i$ and $G$ to be the elements that exist in $L$ after $x_i$. 
Therefore, the cost to access element $x_i$ can be written as follows: $C_{MTF} = |F| + |G| + 1$. 

A key difference betweem this problem and the above is that the initial number of tokens is not necessarily zero, 
since there could exist elements within $G$ since the initial list and $L$ may disagree on element ordering. 
Therefore, we consider the maximum initial number of tokens in $x_i$'s piggy bank. 
Since we consider $1$ token is added based on a pair mismatch, there are $\binom{n}{2}$ pairs in the list, and therefore at most $\binom{n}{2}$ tokens. 

Since $\binom{n}{2} =\frac{n!}{2! \cdot (n-2)!} = \frac{n \cdot (n-1)}{2} = \frac{n^2 - n}{2}$, 
we can set up the following upper bound for the number of tokens: $\text{initial num tokens} \leq \binom{n}{2} < n^2$. 

We consider the the accounting cost $A_i$ for an access to $x_i$ to determine the number of tokens. 
This yields the following equation: $A_i = |F| + |G| + 1 - |G| + |F| = 2|F| + 1$. 

From here, we consider the value for $|F|$. Let us assume that element $x_i$ is in position $k$ in $L$. Therefore, $|F| = k - 1$. 
We can therefore write the accounting cost as $A_i = 2k - 1 < 2k$. 
This establishes a similar inequality as that in the previous problem. To access element $x_i$ in $L$, we pay $k$. 
Therefore, we can set up the following inequality for the access of any one element: $A_{i} < 2k$. 
Summing over all $m$ elements, we get: $\sum_{i}^{m}(A_i) < 2 \cdot C_{static}$

However, we cannot prove the accounting cost to be equal to the actual cost of the MTF operation. 
This is because we begin with a certain number of tokens, which is strictly less than $n^2$. 
To find the actual cost, we must compute the number of tokens $T$ spent over $m$ operations.
We consider $T_0$ and $T_m$ to be the initial and final number of tokens, respectively: $C_{MTF} = \sum_{i}^{m}(A_i) + (T_0 - T_m)$

Since $T_m \geq 0$, we can write this as the following inequality: $C_{MTF} \leq \sum_{i}^{m}(A_i) + T_0$. 
Since the upper bound for the number of tokens is $n^2$, we can substitute that into the equation to yield: $C_{MTF} \leq \sum_{i}^{m}(A_i) + n^2$.

$C_{MTF} < \sum_{i}^{m}(A_i) + n^2$. Therefore, $C_{MTF} - n^2 < \sum_{i}^{m}(A_i) < 2 \cdot C_{static}$. 

Thus, we can conclude: $C_{MTF} < 2 \cdot C_{static} + n^2$.

\end{enumerate}


\section{\texorpdfstring{$k$-th smallest elements (33 points)}{k-th smallest elements (33 points)}}
Given an array $a_1,a_2,\mathellipsis,a_n$ and an integer $k\in[n]$. For each $i\in[n+1-k]$, 
let $b_i$ be the $k$-th smallest element in $a_1,\mathellipsis,a_{k+i-1}$.

\begin{enumerate}
\item (17 points) Design an algorithm which outputs the array $b_1,\mathellipsis,b_{n+1-k}$. The running time should be $O(n\log k)$.

I use a binary heap structure, with a maximum size of $k$ to organize array $A$. We define this heap as $H$ The key of $H$ is the value of each integer in $A$. 
$H$ reverses the typical structure of the binary heap, where instead of organizing in ascending order of the key, it would organize in descending order; 
that is, the root of $H$ would be the largest element, making it a max heap. From here, the first $k$ elements would be inserted into $H$. 
Each insertion would occur at the next open leaf, and then continually swap with its parent while the node is greater than its parents. 
Because during the insertions, each of the new nodes moves into the next open leaf node, where each level of the tree is filled out before adding another, 
the depth of the tree must be less than $(\log x) + 1$ where $x$ is the number of nodes in the tree. Since there are $k$ nodes in the heap, 
the depth of the tree would be $\log k$, meaning each insertion takes at most $O(\log k)$ time. 
This way, after $k$ insertions, the heap contains exactly $k$ elements. Since the root of a max heap is the largest of these $k$ elements, 
it is the $h$th smallest element of the prefix. Thus, a function \verb|find_max()| would be needed to return the $k$th smallest element without extracting it, 
which is simply returning the value of the root, therefore taking $O(1)$ time.

From here, the size of $H$ would remain constant. The algorithm would iterate through each element $i$ of $A$, and check whether $i \geq$ \verb|find_max(H)|. 
If so, the $k$th smallest element remains the same, so no further insertions into the heap must occur. If $i <$ \verb|find_max(H)|, 
then the current maximum must be extracted, and the new element $i$ must be inserted into $H$, to find the new $k$th smallest element. 
The extraction process first involves swapping the rightmost leaf of $H$ with the maximum, removing the maximum, 
and then the new root "swimming down" to the correct location. This involves traversing, at most, the maximum number of levels of the tree, which is $\log k$, 
meaning the extraction process takes at most $O(\log k)$ time. Then, $i$ must be inserted into $H$, which once again, takes $O(\log k)$ time, 
as is proven above. Then, the new root of the tree is the new $k$th smallest element, so \verb|find_max(H)| returns the next $k$th smallest element. 
Therefore, the combination of $2$ $O(\log k)$ and $1$ $O(1)$ operations means each insertion has an upper bound of $O(\log k)$

Generalizing this process, at every point after processing a prefix within $A$ of length at least $k$, 
$H$ contains exactly the $k$ smallest elements of the prefix. Since $H$ is a max heap, its root is the $k$th smallest element of the prefix, 
so it can be extracted to find the $k$th smallest element of any prefix. This repeats for each element of $A$ with a position greater than $k$. 
Thus, for each of the $n$ elements of $A$, an $O(\log k)$ operation (either insertion for the first $k$ elements or 
extraction and insertion for subsequent elements) could occur. Therefore, the overall time complexity needed to form 
an array of the $k$th smallest element of $A$ for every subarray of $A$ which is at least length $k$ is $O(n \log k)$.

\item (16 points) Prove the correctness and the running time for your algorithm.

Correctness Proof:

We assume the correctness of the binary heap structure as per the explanation provided in the lecture slides. 
Therefore, in instantiating a max-heap $H$, we assume that the functions \verb|insert()|, \verb|find_max()|, and \verb|extract_max()| exist, 
and they each perform their respective function described correctly. 
Thus, this proof will center on proving that this binary heap structure is utilized correctly in the algorithm described to solve the listed problem.

I will prove this by induction. I will first define some terms. $A$ is the array we consider. $A$ has $n$ elements, and $A_i$ is the $i$th element of $A$. 
Finally, I define $K(i)$ to be the condition that after processing elements $A[0:i]$, $H$ contains exactly the $k$ smallest elements of $A[0:i]$, 
and \verb|find_max(H)| returns the $k$th smallest element of $A[0:i]$.

Our base case is $i = k-1$. Therefore, the first $k$ elements have been processed and inserted into $H$ correctly, 
so $H$ contains $A[0:(k-1)]$. Since $H$ is a max-heap, then \verb|find_max(H)| correctly returns the root of $H$,, which is by definition, 
the $k$th smallest element of $A[0:(k-1)]$. Therefore, $K(k-1)$ holds.

For our inductive hypothesis, we assume $K(i)$ holds. We then apply the algorithm again. We will prove that this means that $K(i+1)$ will hold true. 
This means that to start, we have max heap $H$ containing the $k$ smallest elements in $A[0:i]$. Now, we examine element $A_{i+1}$. 

This creates two cases:
\begin{enumerate}[label=\arabic*., noitemsep]
    \item $A_{i+1} \geq$ \verb|find_max(H)|
    \item $A_{i+1} <$ \verb|find_max(H)|
\end{enumerate}

The first case is the simplest. If $A_{i+1} \geq$ \verb|find_max(H)|, then the current \verb|find_max(H)| remains the $k$th smallest element of $A[0:(i+1)]$. 
Therefore, the new element is at least as large as the $k$th smallest element. Therefore, the $k$ elements already in $H$ remain the $k$ smallest elements 
in the new prefix $A[0:(i+1)]$, so the $k$th smallest element remains the $k$th smallest element of $A[0:i]$. 
Therefore, in this case, the condition $K(i+1)$ holds. 

We now consider the second case. In this case, $A_{i+1}$ is smaller than the current $k$th smallest element. 
Therefore, $A_{i+1}$ belongs among the $k$ smallest elements of the new prefix, while the previous $k$th smallest element is no longer among those $k$ elements, 
so we must find a new $k$th smallest element. Therefore, we use \verb|extract_max(H)| to pull what was just the $k$th smallest element. 
This creates a tree with $k-1$ nodes. Thus, to find the $k$th smallest element, we must insert $A_{i+1}$ into $H$. Since we know \verb|insert()| works, 
we can assume that after $A_{i+1}$ is inserted into $H$, the root of $H$ is now the biggest element in $H$. 
Because the heap now contains exactly the $k$ smallest elements of $A[0:(i+1)]$, the largest element in the heap is the $k$th smallest element of $A[0:(i+1)]$. 
Thus, performing \verb|find_max(H)| returns what is now the $k$th smallest element. 

Therefore, we have shown that if $K(i)$ is true and the algorithm is applied again, then $K(i+1)$ is true. 
In proving both the base case and the inductive step, we have proven the correctness of the algorithm.

\medbreak

Runtime Proof:

We will first prove that if binary heap $H$ has $k$ elements, it has a depth of at most $(\log k) + 1$. 
In my algorithm, heap $H$ is structured so that each level of the tree must be complete before adding a node on the next level. 
A complete binary tree with $k$ nodes has height $\floor{(\log_2 k)} + 1$, and therefore at most $\floor{(\log_2 k)} + 1$ levels.

We will now prove that \verb|insert()| and \verb|extract_max()| run in $O(\log k)$ time, and \verb|find_max()| runs in $O(1)$ time. 

Starting with the latter, we know that the maximum element of $H$ is the root. Therefore, \verb|find_max(H)| looks at the root of $H$ 
and returns it, regardless of the number of nodes in $H$. Therefore, this process takes $O(1)$ time.

Let us assume we are inserting the $k$th element into $H$. Therefore, the element would slot into the next open node, 
making the tree now have a depth strictly less than $(\log_2 k) + 1$. Putting the node into $H$ is one operation regardless of the number of nodes in $H$, 
so this operation is $O(1)$. However, from here, the new element must find the correct location in the $H$, 
where it is smaller than either of its children and greater than or equal to its parents. This is performed using the "swimming up" idea, 
where the element is swapped with its parent until the above conditions are met. 
Let us consider a case where the element inserted is the new maximum. This element must traverse through every level of the tree. 
Therefore, this process has a time complexity of $O((\log k) + 1)$, which simplifies to $O(\log k)$. 
Generalizing this, we conclude that when inserting into a heap containing at most $k-1$ elements, the new element is placed at the next open position, 
after which it may move upward through at most $O(\log k)$ levels.

Now, we will move on to proving that removing the maximum element of heap $H$, with $k$ elements, has a time complexity of $O(\log k)$. 
First, we must swap the root of $H$ with the last element for simple removal. This takes constant time because 
the root and the last element of the heap are both directly accessible, so they can be swapped in $O(1)$ time. 
Then, the last element is easily deleted, and the tree maintains its complete structure. 
From here, what was previously the last element must "sink down". First, the algorithm checks if it is smaller than either of its children. 
If so, it is swapped with the larger of its children until it either has no children or it is greater than or equal to both its children. 
In the worst case, the element moved to the root is the smallest element in $H$, so it has to move down through every level of the tree.
Therefore, in the sinking down process, it traverses each of the $(\log k) + 1$ levels of the tree, 
meaning that this algorithm has a time complexity of $O((\log k) + 1)$, or $O(\log k)$. 
Finally, the algorithm also spends $O(1)$ time recording each output as the final process. 

Therefore, taken together, the process of extracting the maximum of $H$ with $k$ elements, inserting a new element, and finding the new maximum 
has $2$ $O(\log k)$ and $2$ $O(1)$ processes performed sequentially, meaning the time is $2 \cdot O(\log k) + 2 \cdot O(1)$, 
which simplifies to $O(\log k)$. For each of the $n$ elements of $A$, some combination of these processes occur -- 
either an \verb|insert()| for the first $k$ elements, a \verb|find_max(H)| when the $k$ smallest elements remain the same, or all $3$. 
Thus, since this $O(\log k)$ process occurs $n$ times, the algorithm has an overall time complexity of $O(n \log k)$.

Finallyk, we must consider the case where $k = 1$, since $\log 1 = 0$. When $k = 1$, the heap is of size 1 and keeps track of the $1$st smallest, 
or smallest element of the array so far, and therefore runs in $O(n)$ time. This still maintains an upper bound of $O(n \log k)$, 
so we can conclude that the time complexity of this algorithm for any postive integer $k$ is $O(n \log k)$

\end{enumerate}

\end{document}
